% !PRINTED 431
\weberpage{431}
%<<CONT
\gesperrt{is satisfied. But the same congruence is also satisfied if a number
congruent to $\alpha$ modulo $n-1$ is put as exponent, and in every
complete system of residues one finds only one such number
$\alpha$.}

The exponent $\alpha$ is called the \gesperrt{index} of $a$ with respect to the
\gesperrt{base} $g$, and one writes
\weq{17}{\alpha \equiv \operatorname{ind} a \pmod{n-1}.}

One thus has, for every number not divisible by $n$,
\weq{18}{g^{\operatorname{ind} a} \equiv a \pmod{n}.}

From the two formulae
\weq{19}{g^{\operatorname{ind} a + \operatorname{ind} b} \equiv ab, \quad g^{m \operatorname{ind} a} \equiv a^m \pmod{n}}
there follows the theorem:

6.\ \gesperrt{The index of a product is equal to the sum of the indices
of the factors; the index of the $m$th power of $a$ is equal
to $m$ times the index of $a$.}

These theorems are quite analogous to the corresponding theorems that
relate to computation with logarithms, and the analogy can be pursued
still further. Thus, for example, the transition from one base $g$ to
another base $g'$ is expressed by the relation
\weq{20}{\overset{g}{\operatorname{ind}}\, a \equiv \overset{g}{\operatorname{ind}}\, g' \; \overset{g'}{\operatorname{ind}}\, a \pmod{n-1}}
between the two systems of indices. We shall further adduce the easily proved theorem that among the
numbers $a$ not divisible by $n$, \gesperrt{those and only those} are
primitive roots whose indices are relatively prime to $n-1$.

For practical computations one conveniently makes use of so-called
index tables, which correspond to the tables of logarithms.

In the congruence (17) one calls $\alpha$ the index and $a$ the numerus,
and it is best to set up two tables, of which the one gives for each
index the numerus, the other for each numerus the index, where the
second is obtained from the first by rearrangement. Here it is to be
observed that for the indices the modulus $n-1$, and for the numeri the
modulus $n$, comes into consideration, and congruent numbers count
as equivalent. Thus one obtains, e.g., for $n = 13$ and the base 2
%CONT>>
